\documentclass[journal]{IEEEtran}
\usepackage{amsmath,amssymb,amsfonts,graphicx,booktabs,array,cite,url}
\usepackage[hidelinks]{hyperref}
\hypersetup{pdftitle={Channel Estimation for Rydberg Receivers Using Learned Hankel Filtering},pdfauthor={Abdullah Haatim},pdfsubject={Simulation study of magnitude-only channel estimation}}
\graphicspath{{figures/}}
\interdisplaylinepenalty=2500
\newcommand{\CN}{\mathcal{CN}}
\newcommand{\Hank}{\mathcal{H}}
\DeclareMathOperator{\sigmoid}{sigmoid}
\DeclareMathOperator{\softplus}{softplus}
\DeclareMathOperator{\tr}{tr}
\hyphenation{Ryd-berg mea-sure-ments}
\begin{document}
\title{Channel Estimation for Rydberg Receivers\\Using Learned Hankel Filtering}
\author{Abdullah~Haatim\thanks{A. Haatim is with the Department of Electrical and Electronics Engineering, Birla Institute of Technology and Science, Pilani, India (e-mail: f20220519@pilani.bits-pilani.ac.in).}}
\maketitle
\begin{abstract}
Rydberg receivers use excited atoms to sense radio signals. Magnitude-only
observations do not directly reveal the wireless channel. Channel structure
can help recover it, but fixed filtering may remove useful details. We
study a small learned estimator that alternates measurement fitting with
Hankel filtering. It learns from withheld observations, without true-channel
targets. Simulations show lower average error and processing time than
the tested classical Hankel estimator. Both structure and noise information
help average accuracy. However, the classical method remains better with
a weak reference and in some few-path channels. Learned filtering is a
useful option for the tested model within these operating limits.
\end{abstract}
\begin{IEEEkeywords}
Channel estimation, Hankel filtering, magnitude measurements, Rydberg receivers.
\end{IEEEkeywords}

\section{Introduction}
\IEEEPARstart{R}{ydberg} atoms respond to radio-frequency electric fields,
offering a way to study new wireless receivers. Here, signals from several
users combine with a known reference. The receiver measures their combined
magnitude, not their full complex value. The channel describes changes in
signal strength and phase on the way to the receiver. Model-based methods
and learned methods such as URformer estimate it~\cite{xu2025,xiao2026}.

A few propagation paths can create a repeated pattern across a regular
array. Arranging channel samples in a Hankel matrix lets an estimator use
this pattern. Strong filtering can remove real details; weak filtering can
leave noise. Searching many filtering settings also takes time. We study a
small model that learns how much filtering to apply. The true channel is
never shown to it during training. Instead, some measurements are hidden
from the estimator and used to check its predictions. This is called
self-supervised learning.

Our contributions are: (i) a 188-parameter estimator that combines Rician
measurement updates with smooth Hankel filtering; (ii) training this
estimator on withheld measurements, without channel targets or channel-error
checkpoint selection; and (iii) paired tests of accuracy, processing time
and simpler versions, including fresh-test replication and input-removal
tests. We report weak-reference and few-path failures alongside the gains.

\begin{figure*}[!t]
\centering\includegraphics[width=\textwidth]{fig1_method.pdf}
\caption{Estimator overview. The controller chooses filtering strength and mixing. During training, nine withheld pilots score predictions but do not enter initialization or the ten stages.}
\label{fig:method}
\end{figure*}

\section{Related Work}
Training on withheld measurements is established in image
reconstruction~\cite{yaman2020} and self-supervised denoising is used for
Rydberg sensor signals~\cite{liu2026}. Huang et al.\ learn Hankel-factor
updates from fully sampled targets~\cite{huang2021}. Luo et al.\ combine
self-supervision and structured low-rank reconstruction for complex MRI
data~\cite{luo2023}. DeepTensor fits network-based low-rank factors to an
observed object~\cite{saragadam2022}. Rician-likelihood training for
magnitude data is also established~\cite{parker2023}. Neither
self-supervision, learned Hankel processing nor the Rician loss is new
on its own.

We recover a complex multiuser channel from noisy magnitudes and a known
reference, rather than restoring measured complex samples. A small
controller sets filtering strength instead of generating low-rank factors.
Our companion manuscript studies classical Hankel estimation and rank
selection~\cite{haatimcompanion}; here we learn filtering without channel labels.

Structured alternating least squares~\cite{zachariah2012},
data-driven OptShrink~\cite{nadakuditi2014}, low-rank Hankel denoising
(LRHD)~\cite{yin2021} and robust Hankel completion~\cite{cai2023}
already provide other structural methods. Their original linear or direct
matrix observations differ from our referenced magnitudes. Sparse
amplitude-only beam training is also relevant~\cite{wang2026}, but uses
adaptive sensing without our additive reference. We do not claim an
optimal shrinker or superiority over these complete methods. The
supplement distinguishes their assumptions and reports a limited
post-processing comparison.

\section{Receiver and Channel Model}
An array with $N$ elements receives $P$ known training symbols, or pilots,
from each of $K$ users. The observations are
\begin{equation}
 Z=|GS+B+W|,\qquad W_{np}\sim\CN(0,v).
 \label{eq:model}
\end{equation}
Here $G\in\mathbb C^{N\times K}$ is the channel,
$S\in\mathbb C^{K\times P}$ contains the pilots, and
$B\in\mathbb C^{N\times P}$ is the known reference. Magnitudes act
entrywise. Independent noise has total complex variance $v$, or $v/2$ in
each real component. The reference makes magnitudes depend on relative
signal phases, but does not ensure accurate recovery in every case.
Reference calibration and noise variance are known, except in stated
input-mismatch checks. Receiver conversion gain is absorbed into the scaling.
This effective model does not simulate atomic dynamics or optical readout.

For the half-wavelength uniform linear array,
\begin{equation}
 g_{nk}=\sum_{\ell=1}^{L_k}\alpha_{\ell k}e^{-j\pi n\sin\theta_{\ell k}},
 \quad n=0,\ldots,N-1.
\end{equation}
Independent gains follow $\CN(0,1/L_k)$, giving each user unit average
channel power. Angles are uniform on $[-\pi/2,\pi/2]$; pilot entries
follow $\CN(0,1)$. The reference has a known spatial response and constant waveform. Define
$\mathrm{SNR}=10\log_{10}(K/v)$ and
$\mathrm{RSR}=10\log_{10}|\alpha_b|^2$, where $\alpha_b$ is reference
amplitude. RSR uses one user's power: for three users, RSR 12~dB is
about 7.23~dB relative to their combined power.

\subsection{What the Measurements Can Determine}
For one array element, write $\mu_p=g^Ts_p+b_p$ and
$x=[\operatorname{Re}g^T,\operatorname{Im}g^T]^T$. If $\mu_p\ne0$, the
real Jacobian of $|\mu_p|$ has row
\begin{equation}
 A_p=\left[\operatorname{Re}c_p^T,-\operatorname{Im}c_p^T\right],
 \qquad c_p=\mu_p^*s_p/|\mu_p|.
\end{equation}
Full column rank of $A$ ensures local recovery; a small minimum singular
value means poor stability. Thus $P\ge2K$ is necessary for this condition,
not sufficient. Full row rank of $S$ alone is insufficient. With $B=0$,
channel rotation leaves magnitudes unchanged. A known reference can break
this ambiguity, but its strength alone does not ensure recovery. This
unrestricted-row condition is local, not a global or sparse-channel
guarantee~\cite{gao2018}. Proof and information matrix are in the supplement.

\section{Proposed Estimator}
The estimator starts with a rough channel guess. It repeats two steps ten
times: one pulls the guess toward the measurements, and the other smooths
it using the array structure. A small network decides how much smoothing
to apply at each step (Fig.~\ref{fig:method}).

\subsection{Measurement Update}
The magnitude of a complex signal with Gaussian noise follows a Rician
distribution. For noiseless magnitude $a=|\mu|$, its density is
\begin{equation}
 p(z\mid a,v)=\frac{2z}{v}e^{-(z^2+a^2)/v}I_0(2za/v),
\end{equation}
where $I_0$ is a modified Bessel function. Let $M=GS+B$ and
$\kappa=2Z|M|/v$. Entrywise, form
\begin{equation}
 Q=Z\frac{I_1(\kappa)}{I_0(\kappa)}\frac{M}{|M|},
\end{equation}
using the continuous limit at $M=0$. The expectation-maximization (EM)
update and learned step are
\begin{align}
 G_{\rm EM}&=(Q-B)S^H(SS^H)^{-1},\\
 G_{\rm dc}&=G+\gamma_t(G_{\rm EM}-G),\quad
 \gamma_t=2\sigmoid(a_t).
\end{align}
A linear solve evaluates the inverse. Input pilots must have full row
rank. Ten stages are a fixed computation budget, not a proved optimum.
Separately trained five- and twenty-stage models have not been compared.

\subsection{Smooth Hankel Filtering}
Each updated user channel forms a Hankel matrix $H=\Hank(g_{{\rm dc},k})$,
with equal entries along each antidiagonal. At $N=32$, it is $17\times16$.
Singular values describe component strengths. Hard rank truncation deletes
the smaller components. Our smooth rule is
\begin{equation}
 \widehat\sigma_i=\frac{\sigma_i^3}{\sigma_i^2+\tau^2}.
\end{equation}
Small components shrink more without being deleted outright, but may
contain real channel details.

Set $q=\lceil N/2\rceil$, $s_H=\|H\|_F/\sqrt q$, and
$\tau=s_H[0.001+\softplus(u)]$. The equivalent operation is
\begin{align}
 \widehat H&=H-\tau^2H(H^HH+\tau^2I)^{-1},\\
 g_k^{t+1}&=g_{{\rm dc},k}+\eta_{t,k}
 (\Hank^\dagger\widehat H_k-g_{{\rm dc},k}),
\end{align}
where $\Hank^\dagger$ averages antidiagonals and
$\eta_{t,k}=\sigmoid(w)$. The code normalizes by $s_H$ before solving and
guards zero scales. Gradients pass through the solve, without an
identity-backward approximation. Neither control estimates path count.
The filter does not increase the Hankel Frobenius norm, but this does not imply
smaller channel error. The relaxed EM step decreases its quadratic
surrogate for $0\le\gamma_t\le2$; the following filter need not decrease
the likelihood. We make no convergence claim for the full estimator.

\subsection{Controller Inputs}
The network has eight inputs, 16 hidden units, two outputs and SiLU
activation $x\sigmoid(x)$. Its 178 parameters and ten step sizes give 188. For
$p_i=\sigma_i/\sum_j\sigma_j$, its first
three inputs are
\begin{equation}
 e^{-\sum_i p_i\log p_i}/q,\qquad p_1,\qquad
 \sum_{i=q-\lceil q/2\rceil+1}^{q}p_i.
\end{equation}
They tell the network how spread out the components are, how much the
largest contributes, and how much the smaller components contribute.
Two more inputs measure disagreement with observations and expected
noise relative to estimated channel power:
\begin{align}
 f_4&=\log\left(1+\frac{\operatorname{mean}[(Z-|G_{\rm dc}S+B|)^2]}{v}\right),\\
 f_{5,k}&=\log\left(1+\frac{v\tr[(SS^H)^{-1}]}{K\widehat p_k}\right),
 \quad\widehat p_k=\|g_{{\rm dc},k}\|^2/N.
\end{align}
Inputs $P/(2K)$, $\log q$ and $(t+1)/10$ give measurement count, size and
stage ($t=0,\ldots,9$). Denominator guards avoid division by zero.
True channels, path counts and family labels are not supplied. Weights stay fixed.

\subsection{Training Without True Channels}
Of 30 pilot columns, 21 enter the estimator and nine are withheld.
Training scores predicted withheld magnitudes using the Rician loss
\begin{equation}
 \ell(z,\mu,v)=\frac{(z-|\mu|)^2}{v}-\log I_{0e}(2z|\mu|/v).
\end{equation}
Here $I_{0e}$ is the scaled Bessel function. We omit an observation-only
term and use small-field series at continuous limits. The loss fits
Rician amplitude, not its mean. Splits change each epoch; separate
validation cases with a fixed split select the checkpoint. Deployment
uses all pilots, without extra measurements or weight training.
This avoids test-driven selection; it does not make withheld likelihood
equivalent to channel NMSE or remove dependence on the validation mixture.

\section{Simulation Study}
\subsection{Setup and Comparators}
Table~\ref{tab:setup} gives the settings. Half the training cases have
three to seven paths per user; the rest use fixed counts of 2, 4, 8, 12
and 16. Validation weights 126 SNR, RSR and path-setting combinations
equally. These condition labels set validation weights only; model
selection uses withheld likelihood, not channel error.

\begin{table}[!t]
\caption{Main Simulation and Training Settings}\label{tab:setup}
\centering\small
\begin{tabular}{@{}ll@{}}\toprule
Setting & Value\\\midrule
Array elements / users / pilots & $32/3/30$\\
Training / validation cases & $16{,}384/2{,}048$\\
SNR values (dB) & $-10,-5,0,5,10,15,20$\\
Training RSR values (dB) & $5,12,20$\\
Training seeds & 3, sharing training cases\\
Estimator stages / parameters & $10/188$\\
Training epochs / batch size & $20/32$\\
Optimizer / learning rate & Adam / $0.001$\\
Input / withheld pilots & $21/9$\\
Training / evaluation precision & Single / double\\\bottomrule
\end{tabular}
\end{table}

The main study trains 18 models: two updates, three filtering choices and
three seeds. Seeds change weights, batch order and pilot splits, but share
training cases. Training took 3.67 hours with overlapping workers,
excluding data preparation and preflight. Tests cover 32 conditions with
400 paired cases each. Two noise-input tests reuse cases, giving 12,000
distinct cases. All methods share channels, pilots and noise per case.

HS-GS, the companion study's classical Hankel estimator~\cite{haatimcompanion},
fits ranks $1,\ldots,q$ on 70\% of pilots, using 25 EM updates and four
rank/Hankel projection sweeps per update. Squared magnitude residuals on
the remaining pilots select a rank. It then refits all pilots for 100
updates. Rank search uses no extra observations and counts toward time.
Its rank-selection loss differs from our Rician training loss.
Other controls use EM, Rician or amplitude-gradient updates. Global
filtering has 30 parameters and scales its threshold with $s_H$; no
filtering has ten. Hard-Hankel Rician uses development-selected rank ten
and one projection sweep per update. Constant-output controls replace
adaptive outputs with validation averages without retraining.

Our earlier URformer-style model has 1,586,900 parameters and used 80,000
labelled examples, fixed pilots, $P=20$ and RSR 10~dB. It has no matched
evaluation under this paper's protocol; its earlier scores are not
included in Table~\ref{tab:all}.

Channel error is normalized mean squared error (NMSE):
\begin{equation}
 \mathrm{NMSE}_{\rm dB}=10\log_{10}
 \frac{\sum_i\|\widehat G_i-G_i\|_F^2}{\sum_i\|G_i\|_F^2}.
\end{equation}
Lower is better. Per-seed dB scores and the seven main SNR settings are
averaged equally. Positive gain means lower error. A severe error has
more energy than the true channel (per-case linear NMSE above one).
Pointwise 95\% intervals use 4,000 paired whole-case resamples, conditional
on saved models, excluding training-dataset uncertainty. Classical stopping
methods have a 1,000-update cap; reaching it need not mean convergence.

\begin{figure}[!t]
\centering\includegraphics[width=\columnwidth]{fig2_snr.pdf}
\caption{Main SNR test: $N=32$, $K=3$, $P=30$, RSR 12~dB, and three-to-seven-path channels. Lines average per-seed scores; bands are conditional paired-case 95\% intervals.}
\label{fig:snr}
\end{figure}

\subsection{Accuracy and Processing Time}
\begin{table}[!t]
\caption{All Main-Study Comparators}\label{tab:all}
\centering\footnotesize\setlength{\tabcolsep}{3pt}
\begin{tabular}{@{}lrrr@{}}\toprule
Method & NMSE & Time & Severe\\
 & (dB) & (ms) & (\%)\\\midrule
\csname @@input\endcsname baseline_rows.tex
\bottomrule\end{tabular}
\parbox{\columnwidth}{\footnotesize NMSE and severe-error rates average seven primary SNR conditions equally. Time is the median at SNR 5~dB, RSR 12~dB. Learned scores average three seeds; classical cases are not counted three times.}
\end{table}

HS-GS is the strongest original classical comparator on the declared average
(Table~\ref{tab:all}). The adaptive gradient-descent (GD) core has lower
test NMSE than the primary EM core. We retain EM because validation
likelihood selected it before test generation.

The EM model's average gain over HS-GS is 1.675~dB, with interval
$[1.618,1.730]$~dB and seed-specific gains of 1.533--1.787~dB
(Fig.~\ref{fig:snr}). At 20~dB, the gain is 0.139~dB with interval
$[-0.017,0.294]$~dB, which does not establish a clear improvement.

Severe errors fall from 14.25\% for HS-GS to 0.45\% for the proposed
method. In the tested SNR mixture, that is roughly one badly wrong
estimate in seven cases versus one in 220.

Median time is 8.24~ms versus 525.62~ms for HS-GS, about 64 times faster
(Table~\ref{tab:all}). Timing includes initialization and rank search,
using one Apple M2 CPU thread, double precision, batch size one,
100 cases and three repeats at SNR 5~dB and RSR 12~dB.
These are implementation-specific measurements:
learned methods use PyTorch and classical methods use NumPy/SciPy.
Offline training is separate. No filtering is faster but less accurate.
With $m=N-q+1$, one adaptive stage costs
$O(NKP+NK^2+K^3+Kmq^2+Kq^3)$: the spectral features still require an SVD.
HS-GS repeats comparable projections across candidate ranks and updates.
The supplement separates initialization, update and projection costs;
188 weights do not imply negligible computation.

\begin{figure*}[!t]
\centering\includegraphics[width=\textwidth]{fig3_limits.pdf}
\caption{Limits at SNR 5~dB: (a) reference strength with three-to-seven-path channels; (b) fixed path counts at RSR 12~dB. Positive gain favors the proposed estimator. Bars are conditional pointwise paired-case 95\% intervals. Path counts overlap training.}
\label{fig:limits}
\end{figure*}

\subsection{Operating Limits}
The proposed method loses 1.617~dB at RSR 0~dB, 1.104~dB with two paths,
and 0.203~dB with three paths (Fig.~\ref{fig:limits}). Good average
performance does not ensure good performance with a weak reference or
a simple channel.

Transfer tests give gains for noise inputs wrong by $\pm3$~dB and three
channel alternatives: four clusters of ten rays within $\pm5^\circ$,
ten equal-power paths, or ten paths with geometrically decreasing powers.
Alternatives match base-channel energy and are tested only at SNR 5~dB,
RSR 12~dB. Dimension tests use $N\in\{8,16,32,64\}$,
$K\in\{2,3,4\}$ and $P\in\{10,20,30\}$, changing one at a time.
At $N=16$ and $64$, gains are 1.223 and 0.165~dB; at $N=8$ the
interval includes zero. No tested pilot budget approaches $2K$ closely.
The supplement reports all dimension changes. Path counts
overlap training; these tests do not resolve failures under earlier protocols.

\subsection{Further Comparisons}
A follow-up trains 15 models and uses two fresh test sets, each with
seven SNR settings and two difficult settings of 200 cases.
Six primary comparisons are adjusted for multiple testing per set.
Independent training-data runs gain 1.939 and 1.837~dB over HS-GS;
each dataset uses one seed, limiting replication. These are separate
from the main 1.675~dB result. The supplement retains all six comparisons.

A matched-architecture supervised recipe uses mean linear NMSE and the
same input pilots and training budget. Its approximately 1~dB gap depends
on checkpoint and averaging choices: it reduces to about 0.73 or 0.36~dB
under alternative rules. Supervision wins in some settings, including two
paths. This is not a general superiority claim for self-supervision.

Repeated filtering gains about 0.70~dB over separately trained final-only
filtering. A separate timing sample gives 4.65 versus 2.05~ms. Mixing is
usually near one; fixing it at one gives a small average gain but worsens
two of three seeds on two-path cases. The controller adjusts shrinkage
strength, not a reliable on/off decision.

\subsection{Which Inputs Help?}
We compare full, spectrum-only, noise/residual-only and size/stage-only
controllers. Size and stage remain in every version. All use the same
update and filter, trained from scratch on 4,096 shared cases, with three
seeds and 12 epochs. Two fresh sets contain 2,304 cases: seven SNR
settings and two difficult settings, 128 cases per condition per set.

Full inputs gain 1.22--1.25~dB over spectrum-only, 0.57--0.58~dB over
noise/residual-only, and 2.21--2.23~dB over size/stage-only inputs;
adjusted intervals remain positive. Yet they lose 0.39--0.40~dB against
noise/residual-only at SNR 5~dB, failing the declared 0.25~dB protection
requirement, and 0.82--0.93~dB on two-path cases. The supplement retains
the figure. Both groups help the average, not every channel. All variants
store 188 parameters with masked weights inactive; noise variance still
enters the measurement update and channel information still enters the filter.

\subsection{Small Reviewer Diagnostics}
Fresh checks use 32 cases per cell and frozen weights, not a new
confirmation study. Rician rank scoring improves HS-GS by 0.182~dB;
the learned model retains a 1.282~dB average gain. A separately
development-tuned EM100-plus-OptShrink control narrows the learned gain
to 0.552~dB and wins at SNR 20~dB and RSR 0~dB. In matched one-pass
post-processing, OptShrink beats cubic filtering by 2.042~dB. Thus the
cubic rule itself is not superior. Supplying noise variance 6~dB too
high also makes the learned model 1.409~dB worse than Rician-ranked
HS-GS. The supplement reports paired intervals, reference-input errors
and controller trajectories. These results strengthen the comparisons
but narrow the useful operating range.

\section{Discussion and Conclusion}
A small model can learn useful Hankel filtering without true-channel
targets. Its practical setting is repeated estimation near the simulated
training population, where offline training is reused. The average
accuracy and time gains come with weak-reference and few-path losses.

All results are simulations. Reference calibration is assumed known,
and no physical receiver measurements are used. Three seeds sharing
training cases are not three independent datasets. We have not proved
unique channel recovery in every setting or convergence of the full
learned estimator. The comparisons also do not establish superiority
over published URformer performance.
Noise-dependent regularized classical controls and broader independent
training/pilot replications would be needed for a stronger robustness
claim. We do not infer the cause of a loss solely from its association
with SNR, path count or reference strength.
Two-path and RSR 5~dB examples already occur in training, but RSR 0~dB
does not. Changing
their weights or adding a curriculum has not been tested. No robustness
to physical shot statistics, unknown reference drift or two-dimensional
arrays is established. A block-Hankel extension would require a separate
geometry model and validation, not just reshaping the input.

The supplement records initialization, numerical safeguards, controller
trajectories and bounded reviewer diagnostics. These exploratory checks
do not replace the original frozen confirmation results. The original
reproducibility package contains compact outputs and a replay subset,
not every observation/prediction array.
\bibliographystyle{IEEEtran}
\bibliography{references}
\end{document}
